Se houvesse semelhanca geometrica estrita entre robin e stilt, poderiamos usar a mesma lei de escala do Problema 3.30:
\[
l\propto W^{1/3}.
\]
Com $W_s=4.5$ oz, $l_s=8$ in e $W_r=2$ oz,
\[
\frac{l_r}{8}=\left(\frac{2}{4.5}\right)^{1/3}\approx 0.763,
\]
de modo que
\[
l_r\approx 8(0.763)\approx 6.1\ \text{in}.
\]

Esse valor, no entanto, nao e confiavel biologicamente. O robin nao tem a mesma forma corporal nem a mesma funcao locomotora do stilt: suas pernas sao proporcionalmente menores e seu corpo nao e uma copia em escala. Portanto, a regra de semelhanca pode ser aplicada como exercicio de escala, mas nao como previsao realista para a especie.

